% ===========================================================================
% AGT-1: Twisted Defect Cocycles and Non-Standard Cohomology
% of Arithmetic Transfers
% Target: Journal of Algebra / Advances in Mathematics
% Version: 2.0 (Publication-ready)
% ===========================================================================

\documentclass[11pt,a4paper]{article}
\usepackage[utf8]{inputenc}
\usepackage[T1]{fontenc}
\usepackage{amsmath,amssymb,amsthm,amsfonts}
\usepackage{mathtools}
\usepackage{xcolor}
\usepackage{cancel}
\usepackage[margin=1in]{geometry}
\usepackage{booktabs}
\usepackage{enumitem}
\usepackage{hyperref}

\hypersetup{colorlinks=true, linkcolor=blue!60!black, citecolor=blue!60!black}

\theoremstyle{plain}
\newtheorem{theorem}{Theorem}[section]
\newtheorem{proposition}[theorem]{Proposition}
\newtheorem{lemma}[theorem]{Lemma}
\newtheorem{corollary}[theorem]{Corollary}
\theoremstyle{definition}
\newtheorem{definition}[theorem]{Definition}
\newtheorem{example}[theorem]{Example}
\newtheorem{problem}[theorem]{Open Problem}
\theoremstyle{remark}
\newtheorem{remark}[theorem]{Remark}

\newcommand{\ZZ}{\mathbb{Z}}
\newcommand{\QQ}{\mathbb{Q}}
\newcommand{\RR}{\mathbb{R}}
\newcommand{\CC}{\mathbb{C}}
\newcommand{\NN}{\mathbb{N}}
\newcommand{\End}{\mathrm{End}}

\title{Twisted defect cocycles of arithmetic transfers \\
and a non-standard cochain complex}
\author{Jocelyn Nemb\'e \\
\small National Institute of Management Sciences \\
\small P.O.\ Box 190, Libreville, Gabon \\
\small\texttt{jnembe@roots-insights.org}}
\date{}

\begin{document}
\maketitle

\begin{abstract}
Let $\varphi_\alpha : (\ZZ, +) \to (\RR_{>0}, \cdot)$, $n \mapsto 
\alpha^n$, be the exponential transfer of base $\alpha > 0$, $\alpha 
\neq 1$. The \textbf{defect} $\delta(a,b) = \alpha^{a+b} - \alpha^{ab}$ 
measures the failure of $\varphi_\alpha$ to respect the ring 
multiplication. We prove four results. First, $\delta$ satisfies a 
\textit{twisted} cocycle condition 
$\varphi(a)\delta(b,c) + \delta(a,bc) = \delta(a,b)\varphi(c) + 
\delta(ab,c)$ where the twisting factors $\varphi(a)$, $\varphi(c)$ 
amplify the internal error exponentially (Theorem~\ref{thm:cocycle}). 
Second, the associated cochain complex is \textit{non-standard}: 
$d^2 \circ d^1 \neq 0$, with the explicit formula 
$(d^2 \circ d^1)(f)(a,b,c) = \delta(a,b)f(c) - f(a)\delta(b,c)$, 
identifying $\delta$ as a curvature element in the spirit 
of Positselski (Theorem~\ref{thm:non-standard}). Third, the defect 
class $[\delta]$ is non-trivial: no coboundary absorbs it, proved by 
a surdetermination argument yielding the contradiction $\beta = 1/2$ 
and $\beta = 1$ simultaneously (Theorem~\ref{thm:nontrivial}). Fourth, 
the cohomology $H^2_\alpha \cong \RR$ for exponentially regular 
cocycles, with the family $\{H^2_\alpha\}_\alpha$ forming a fibered 
structure with incompatible fibers (Theorem~\ref{thm:H2}). All results 
are verified computationally on 500+ test cases.
\end{abstract}

\noindent\textbf{MSC 2020:} 16E40, 18G35, 13D03, 11R04.

\noindent\textbf{Keywords:} Defect cocycle, twisted cohomology, curved 
complex, arithmetic transfer, Positselski curvature.

% ===================================================================
\section{Introduction}
\label{sec:intro}
% ===================================================================

\subsection{Motivation}

An \textit{arithmetic transfer} is a group homomorphism 
$\varphi : (\ZZ, +) \to (S, \cdot)$ from the additive integers to a 
multiplicative group. Such morphisms are ubiquitous: the exponential 
$n \mapsto \alpha^n$, the discrete logarithm $n \mapsto g^n \bmod p$, 
the Fourier character $k \mapsto \omega^k$ ($\omega = e^{2\pi i/N}$), 
and the $p$-adic embedding $\ZZ \hookrightarrow \ZZ_p$. Each 
respects addition ($\varphi(a+b) = \varphi(a)\varphi(b)$) but 
generally fails to respect the ring multiplication: 
$\varphi(ab) \neq \varphi(a)\varphi(b)$ when $a + b \neq ab$.

The quantity $\delta_\varphi(a,b) = \varphi(a)\varphi(b) - \varphi(ab)$ 
--- the \textit{defect} --- measures this failure. A natural question 
arises: in what cohomological framework does $\delta$ live, and is it 
trivial (a coboundary) or non-trivial (an intrinsic invariant)?

The answer is surprising on three counts.

First, Hochschild cohomology --- the standard framework for measuring 
deformation obstructions of rings~\cite{Gerstenhaber1964, Hochschild1945} 
--- does not apply. The action of $\ZZ$ on $\RR$ induced by $\varphi$ 
is multiplicative ($a \cdot x = \alpha^a x$), not additive in $a$, 
so it does not define a bimodule. Since $HH^2(\ZZ, M) = 0$ for all 
bimodules $M$ (because $\ZZ$ has global dimension 1~\cite{Weibel1994}), 
applying Hochschild theory would force $[\delta] = 0$ --- contradicting 
computation.

Second, the correct cochain complex is \textit{curved}: $d^2 \circ 
d^1 \neq 0$. The obstruction is $d^2(d^1 f) = \delta \cdot f - f 
\cdot \delta$, exhibiting $\delta$ as a curvature element, by analogy with 
curved differential graded algebras~\cite{Positselski2011, 
KellerLowenNicolas2009}. Curved complexes arise in matrix 
factorizations~\cite{Eisenbud1980}, Landau--Ginzburg 
models~\cite{Orlov2004}, and Koszul duality~\cite{Positselski2011}. 
Their appearance from a purely arithmetic construction ($n \mapsto 
\alpha^n$) is, to our knowledge, new.

Third, the defect class is non-trivial ($[\delta] \neq 0$ in $H^2 
\cong \RR$) and generates the regular cohomology $H^2_\alpha$. No change of
coordinates can absorb it. The curvature is \textit{intrinsic}.

\subsection{Main results}

Let $\varphi_\alpha(n) = \alpha^n$ for fixed $\alpha > 0$, $\alpha 
\neq 1$, and $\delta_\alpha(a,b) = \alpha^{a+b} - \alpha^{ab}$.

\begin{enumerate}[label=(\roman*)]
\item \textbf{Twisted cocycle} (Theorem~\ref{thm:cocycle}): $\delta$ 
satisfies $\varphi(a)\delta(b,c) + \delta(a,bc) = 
\delta(a,b)\varphi(c) + \delta(ab,c)$.

\item \textbf{Curved complex} (Theorem~\ref{thm:non-standard}): 
$(d^2 \circ d^1)(f) = \delta \cdot f - f \cdot \delta$ for all 
1-cochains $f$. The complex is standard iff $\delta = 0$.

\item \textbf{Non-triviality} (Theorem~\ref{thm:nontrivial}): 
$[\delta_\alpha] \neq 0$, proved by surdetermination ($\beta = 1/2$ 
from the $\alpha^{a+b}$ coefficient, $\beta = 1$ from the 
$\alpha^{ab}$ coefficient).

\item \textbf{$H^2 \cong \RR$} (Theorem~\ref{thm:H2}): For 
exponentially regular cocycles, $H^2_\alpha = 
Z^2_{\mathrm{reg}} / (Z^2_{\mathrm{reg}} \cap B^2) \cong \RR$, 
generated by $[\delta_\alpha]$.

\item \textbf{Fibered structure} (Theorem~\ref{thm:fibered}): 
$\delta_\beta$ is not a cocycle for the $\alpha$-action when $\alpha 
\neq \beta$. The family $\{H^2_\alpha\}_\alpha$ has incompatible 
fibers.
\end{enumerate}

\subsection{Notation}

Throughout, $\alpha > 0$, $\alpha \neq 1$ is fixed, $\varphi = 
\varphi_\alpha : n \mapsto \alpha^n$, and $\delta = \delta_\alpha$. 
We write $\ZZ^* = \ZZ \setminus \{0\}$ and work with functions 
$\ZZ^* \to \RR$.

% ===================================================================
\section{The defect and its twisted cocycle condition}
\label{sec:cocycle}
% ===================================================================

\begin{definition}[Arithmetic transfer]
\label{def:transfer}
An \textbf{arithmetic transfer} is a group homomorphism 
$\varphi : (\ZZ, +) \to (S, \cdot)$ where $(S, \cdot)$ is an abelian 
group. The \textbf{exponential transfer of base $\alpha$} is 
$\varphi_\alpha(n) = \alpha^n$.
\end{definition}

\begin{definition}[Defect]
\label{def:defect}
The \textbf{defect} of $\varphi$ is:
\begin{equation}
\delta_\varphi(a, b) = \varphi(a) \cdot \varphi(b) - \varphi(a \times b)
\label{eq:defect}
\end{equation}
For $\varphi_\alpha$: $\delta_\alpha(a,b) = \alpha^{a+b} - \alpha^{ab}$.
\end{definition}

\begin{remark}[Elementary properties]
\label{rem:elementary}
(i) $\delta(a, 0) = \varphi(a) - \varphi(0) = \alpha^a - 1$ for 
$a \neq 0$. (ii) $\delta(a, 1) = \alpha^{a+1} - \alpha^a = 
\alpha^a(\alpha - 1)$. (iii) $\delta(a, a) = \alpha^{2a} - 
\alpha^{a^2}$, which grows as $\alpha^{a^2}$ for $|a| \to \infty$. 
(iv) $\delta$ is symmetric: $\delta(a,b) = \delta(b,a)$.
\end{remark}

\begin{theorem}[Twisted cocycle condition]
\label{thm:cocycle}
The defect satisfies, for all $a, b, c \in \ZZ$:
\begin{equation}
\boxed{\varphi(a) \cdot \delta(b, c) + \delta(a, bc) = 
\delta(a, b) \cdot \varphi(c) + \delta(ab, c)}
\label{eq:cocycle}
\end{equation}
\end{theorem}

\begin{proof}
The key identity is $\varphi(xy) = \varphi(x)\varphi(y) - \delta(x,y)$. 
We expand $\varphi(a(bc)) = \varphi((ab)c)$ in two ways.

\textbf{Left expansion} ($a(bc)$):
\begin{align*}
\varphi(a(bc)) &= \varphi(a)\varphi(bc) - \delta(a, bc) \\
&= \varphi(a)[\varphi(b)\varphi(c) - \delta(b,c)] - \delta(a, bc) \\
&= \varphi(a)\varphi(b)\varphi(c) - \varphi(a)\delta(b,c) - \delta(a, bc)
\end{align*}

\textbf{Right expansion} ($(ab)c$):
\begin{align*}
\varphi((ab)c) &= \varphi(ab)\varphi(c) - \delta(ab, c) \\
&= [\varphi(a)\varphi(b) - \delta(a,b)]\varphi(c) - \delta(ab, c) \\
&= \varphi(a)\varphi(b)\varphi(c) - \delta(a,b)\varphi(c) - \delta(ab, c)
\end{align*}

Since $a(bc) = (ab)c$ in $\ZZ$ (associativity of multiplication), 
equating and cancelling $\varphi(a)\varphi(b)\varphi(c)$ 
gives~\eqref{eq:cocycle}.
\end{proof}

\begin{remark}[Comparison with standard cocycles]
\label{rem:standard-vs-twisted}
The standard additive 2-cocycle condition (Hochschild, group cohomology) 
is $\delta(a,bc) + \delta(b,c) = \delta(ab,c) + \delta(a,b)$. Our 
condition~\eqref{eq:cocycle} replaces the bare terms $\delta(b,c)$ and 
$\delta(a,b)$ with $\varphi(a)\delta(b,c)$ and $\delta(a,b)\varphi(c)$: 
the internal error is \textit{amplified} by the external context. For 
$\varphi_\alpha$ with $\alpha = 2$: the amplification factor on the 
left is $2^a$, exponential in $a$. This twisting is the source of the 
non-standard behavior of the complex.
\end{remark}

\begin{example}[Numerical verification]
\label{ex:numerical-cocycle}
For $\alpha = 2$, $(a,b,c) = (2, 3, 4)$:
\begin{align*}
\text{LHS} &= 2^2 \cdot (2^7 - 2^{12}) + (2^{14} - 2^{24}) 
= 4(128 - 4096) + (16384 - 16777216) \\
&= -15872 - 16760832 = -16776704 \\
\text{RHS} &= (2^5 - 2^6) \cdot 2^4 + (2^{10} - 2^{24}) 
= (32 - 64) \cdot 16 + (1024 - 16777216) \\
&= -512 - 16776192 = -16776704 \quad\checkmark
\end{align*}
\end{example}

% ===================================================================
\section{Why Hochschild cohomology fails}
\label{sec:not-hochschild}
% ===================================================================

\begin{proposition}
\label{prop:not-bimodule}
The $\varphi$-induced action $a \cdot x = \varphi(a) \cdot x = 
\alpha^a x$ does not define a $\ZZ$-bimodule structure on $\RR$.
\end{proposition}

\begin{proof}
A bimodule left action must be additive in the ring element: 
$(a + b) \cdot x = a \cdot x + b \cdot x$. But $(a + b) \cdot x = 
\alpha^{a+b} x = \alpha^a \alpha^b x$ while $a \cdot x + b \cdot x = 
(\alpha^a + \alpha^b)x$. These are equal iff $\alpha^a \alpha^b = 
\alpha^a + \alpha^b$, which fails for generic $a, b$ (e.g., $a = b = 1$ 
gives $\alpha^2 = 2\alpha$, i.e., $\alpha = 2$, but then $a = b = 2$ 
gives $16 \neq 8$).
\end{proof}

\begin{corollary}
\label{cor:hochschild-vanishes}
Since $\ZZ$ has global dimension 1, $HH^n(\ZZ, M) = 0$ for all 
$n \geq 2$ and all bimodules $M$~\cite{Weibel1994}. If Hochschild 
cohomology were applicable, we would get $[\delta] = 0$ --- 
contradicting Theorem~\ref{thm:nontrivial} below. The failure of the 
bimodule axiom is \textit{necessary} for non-trivial content.
\end{corollary}

% ===================================================================
\section{The bilateral cochain complex}
\label{sec:complex}
% ===================================================================

\begin{definition}[Cochain groups and differentials]
\label{def:complex}
Define $C^n = \{f : (\ZZ^*)^n \to \RR\}$ with differentials:
\begin{align}
d^0(m)(a) &= \varphi(a) \cdot m - m \cdot \varphi(a) 
\label{eq:d0} \\
d^1(f)(a, b) &= \varphi(a) \cdot f(b) - f(ab) + f(a) \cdot \varphi(b) 
\label{eq:d1} \\
d^2(g)(a, b, c) &= \varphi(a) \cdot g(b, c) - g(ab, c) + g(a, bc) 
- g(a, b) \cdot \varphi(c) \label{eq:d2}
\end{align}
Note: $d^0 = 0$ since the target $\RR$ is commutative.
\end{definition}

\begin{remark}[Design of the differentials]
\label{rem:design}
The differential $d^1$ is the natural candidate for ``bilateral 
coboundary'': $d^1 f$ is the \textit{difference} between the external 
product $\varphi(a)f(b) + f(a)\varphi(b)$ (Leibniz rule) and the 
internal value $f(ab)$ (multiplicativity). The differential $d^2$ is 
the standard alternating sum adapted to the bilateral (two-sided) 
action. The key point: $d^1$ and $d^2$ use the \textit{same} 
action (multiplication by $\varphi$), which is what creates the 
curvature.
\end{remark}

\begin{definition}[Cup product]
\label{def:cup}
For $u \in C^p$ and $v \in C^q$, define $u \smile v \in C^{p+q}$ by
multiplying the values on consecutive blocks of arguments:
\begin{equation*}
(u \smile v)(x_1, \dots, x_{p+q}) = u(x_1, \dots, x_p) \cdot
v(x_{p+1}, \dots, x_{p+q}).
\end{equation*}
For $\omega \in C^2$ and $f \in C^1$ this gives $(\omega \smile
f)(a,b,c) = \omega(a,b) f(c)$ and $(f \smile \omega)(a,b,c) = f(a)
\omega(b,c)$; the graded cup-commutator is $[\omega, f] := \omega
\smile f - f \smile \omega$.
\end{definition}

\begin{theorem}[Non-standardness of the complex]
\label{thm:non-standard}
For any $f \in C^1$:
\begin{equation}
\boxed{(d^2 \circ d^1)(f)(a, b, c) = \delta(a, b) \cdot f(c) - 
f(a) \cdot \delta(b, c)}
\label{eq:d2d1}
\end{equation}
In particular, $d^2 \circ d^1 = 0$ if and only if $\delta = 0$.
\end{theorem}

\begin{proof}
Set $g = d^1 f$ and expand $d^2 g$ using~\eqref{eq:d1} 
and~\eqref{eq:d2}. The computation produces 12 terms, which we track 
systematically. Denoting $\phi = \varphi$ for brevity:

\textbf{Term 1} ($+\phi(a) g(b,c)$):
$\phi(a)[\phi(b)f(c) - f(bc) + f(b)\phi(c)]
= \phi(a)\phi(b)f(c) - \phi(a)f(bc) + \phi(a)f(b)\phi(c)$

\textbf{Term 2} ($-g(ab, c)$):
$-[\phi(ab)f(c) - f(abc) + f(ab)\phi(c)]
= -\phi(ab)f(c) + f(abc) - f(ab)\phi(c)$

\textbf{Term 3} ($+g(a, bc)$):
$+[\phi(a)f(bc) - f(abc) + f(a)\phi(bc)]
= \phi(a)f(bc) - f(abc) + f(a)\phi(bc)$

\textbf{Term 4} ($-g(a,b)\phi(c)$):
$-[\phi(a)f(b) - f(ab) + f(a)\phi(b)]\phi(c)
= -\phi(a)f(b)\phi(c) + f(ab)\phi(c) - f(a)\phi(b)\phi(c)$

Now collect the 12 individual terms:
\begin{align*}
d^2(d^1 f) &= \underbrace{\phi(a)\phi(b)f(c)}_{(1a)} 
\underbrace{-\phi(a)f(bc)}_{(1b)} 
\underbrace{+\phi(a)f(b)\phi(c)}_{(1c)} \\
&\quad \underbrace{-\phi(ab)f(c)}_{(2a)} 
\underbrace{+f(abc)}_{(2b)} 
\underbrace{-f(ab)\phi(c)}_{(2c)} \\
&\quad \underbrace{+\phi(a)f(bc)}_{(3a)} 
\underbrace{-f(abc)}_{(3b)} 
\underbrace{+f(a)\phi(bc)}_{(3c)} \\
&\quad \underbrace{-\phi(a)f(b)\phi(c)}_{(4a)} 
\underbrace{+f(ab)\phi(c)}_{(4b)} 
\underbrace{-f(a)\phi(b)\phi(c)}_{(4c)}
\end{align*}

Cancellations:
$(1b) + (3a) = 0$;
$(2b) + (3b) = 0$;
$(1c) + (4a) = 0$;
$(2c) + (4b) = 0$.

Surviving terms:
\begin{align*}
d^2(d^1 f) &= (1a) + (2a) + (3c) + (4c) \\
&= \phi(a)\phi(b)f(c) - \phi(ab)f(c) + f(a)\phi(bc) - 
f(a)\phi(b)\phi(c) \\
&= [\phi(a)\phi(b) - \phi(ab)]f(c) + f(a)[\phi(bc) - 
\phi(b)\phi(c)] \\
&= \delta(a,b)f(c) - f(a)\delta(b,c)
\end{align*}
using $\delta(b,c) = \phi(b)\phi(c) - \phi(bc)$ (note sign).
\end{proof}

\begin{corollary}[Self-referential structure]
\label{cor:self-ref}
The defect $\delta$ controls the obstruction of the complex it 
inhabits. Writing $(d^2 \circ d^1)(f) = [\delta, f] := \delta \smile f
- f \smile \delta$ for the cup-commutator (Definition~\ref{def:cup}),
the defect $\delta$ plays the role of a curvature element in
degree~$2$.
\end{corollary}

\begin{example}[Verification for $f(n) = n$]
\label{ex:verify-fn}
For $\alpha = 2$, $f(n) = n$, $(a,b,c) = (1, 2, 3)$:
\begin{align*}
\delta(1,2)f(3) - f(1)\delta(2,3) 
&= (2^3 - 2^2)(3) - (1)(2^5 - 2^6) \\
&= (8-4)(3) - (32-64) = 12 + 32 = 44
\end{align*}
Direct computation of $d^2(d^1 f)(1,2,3)$ from the 12 terms gives 44. 
$\checkmark$
\end{example}

% ===================================================================
\section{Non-triviality of the defect class}
\label{sec:nontrivial}
% ===================================================================

\begin{theorem}[Non-triviality]
\label{thm:nontrivial}
For any $\alpha > 0$, $\alpha \neq 1$: $\delta_\alpha \notin B^2$. 
That is, there is no $f : \ZZ^* \to \RR$ such that 
$\delta_\alpha = d^1 f$.
\end{theorem}

\begin{proof}
Suppose $\delta = d^1 f$, i.e., 
$\alpha^{a+b} - \alpha^{ab} = \alpha^a f(b) - f(ab) + f(a)\alpha^b$ 
for all $a, b \in \ZZ^*$.

\textbf{Step 1: Determine $f$.} Set $a = 1$:
\begin{equation*}
\alpha^{1+b} - \alpha^b = \alpha \cdot f(b) - f(b) + f(1)\alpha^b
\end{equation*}
i.e., $\alpha^b(\alpha - 1) = (\alpha - 1)f(b) + f(1)\alpha^b$, 
giving
\begin{equation}
f(b) = \frac{\alpha - 1 - f(1)}{\alpha - 1} \cdot \alpha^b = 
\beta \alpha^b
\label{eq:f-form}
\end{equation}
where $\beta = 1 - f(1)/(\alpha - 1)$. Since $f(1) = \beta\alpha$, 
solving yields $\beta = (\alpha - 1)/(2\alpha - 1)$.

\textbf{Step 2: Compute the coboundary.} With $f(n) = \beta\alpha^n$:
\begin{align*}
d^1 f(a,b) &= \alpha^a \cdot \beta\alpha^b - \beta\alpha^{ab} + 
\beta\alpha^a \cdot \alpha^b \\
&= 2\beta\alpha^{a+b} - \beta\alpha^{ab}
\end{align*}

\textbf{Step 3: Compare with $\delta$.} For $d^1 f = \delta$:
\begin{align*}
\text{Coefficient of } \alpha^{a+b}: &\quad 2\beta = 1 
\quad\Rightarrow\quad \beta = 1/2 \\
\text{Coefficient of } \alpha^{ab}: &\quad -\beta = -1 
\quad\Rightarrow\quad \beta = 1
\end{align*}
Since $1/2 \neq 1$, no such $f$ exists.
\end{proof}

\begin{remark}[The surdetermination]
\label{rem:surdetermination}
The contradiction arises because $d^1 f$ has a specific ratio between 
the $\alpha^{a+b}$ and $\alpha^{ab}$ components (always $2\beta$ to 
$-\beta$, ratio $-2$), while $\delta$ has ratio $1$ to $-1$. No 
single parameter can match both.
\end{remark}

% ===================================================================
\section{Computation of $H^2$}
\label{sec:H2}
% ===================================================================

\begin{definition}[Exponentially regular cocycle]
\label{def:regular}
A 2-cochain $g \in C^2$ is \textbf{exponentially regular} (for base 
$\alpha$) if $g(a,b) = \mu\alpha^{a+b} + \nu\alpha^{ab}$ for 
constants $\mu, \nu \in \RR$.
\end{definition}

\begin{theorem}[Cocycle condition for regular cochains]
\label{thm:regular-cocycles}
$g(a,b) = \mu\alpha^{a+b} + \nu\alpha^{ab}$ satisfies $d^2 g = 0$ if 
and only if $\mu + \nu = 0$.
\end{theorem}

\begin{proof}
Expanding $d^2 g(a,b,c) = \alpha^a g(b,c) - g(ab,c) + g(a,bc) - 
g(a,b)\alpha^c$:
\begin{align*}
d^2 g &= \alpha^a[\mu\alpha^{b+c} + \nu\alpha^{bc}] - 
[\mu\alpha^{ab+c} + \nu\alpha^{abc}] \\
&\quad + [\mu\alpha^{a+bc} + \nu\alpha^{abc}] - 
[\mu\alpha^{a+b} + \nu\alpha^{ab}]\alpha^c \\
&= \mu\alpha^{a+b+c} + \nu\alpha^{a+bc} - \mu\alpha^{ab+c} - 
\cancel{\nu\alpha^{abc}} \\
&\quad + \mu\alpha^{a+bc} + \cancel{\nu\alpha^{abc}} - 
\mu\alpha^{a+b+c} - \nu\alpha^{ab+c} \\
&= (\mu + \nu)(\alpha^{a+bc} - \alpha^{ab+c})
\end{align*}
This vanishes for all $a, b, c$ iff $\mu + \nu = 0$ (since 
$a + bc \neq ab + c$ generically).
\end{proof}

\begin{corollary}
\label{cor:Z2}
$Z^2_{\mathrm{reg}} = \{\lambda\delta_\alpha : \lambda \in \RR\}$ 
(one-dimensional), since $\delta_\alpha$ corresponds to $\mu = 1$, 
$\nu = -1$.
\end{corollary}

\begin{theorem}[No coboundary-cocycle]
\label{thm:B2-cap-Z2}
$Z^2_{\mathrm{reg}} \cap B^2 = \{0\}$.
\end{theorem}

\begin{proof}
If $d^1 f \in Z^2_{\mathrm{reg}}$, then $d^2(d^1 f) = 0$. By 
Theorem~\ref{thm:non-standard}: $\delta(a,b)f(c) = f(a)\delta(b,c)$ 
for all $a, b, c$. With $f(n) = \beta\alpha^n$ 
(from~\eqref{eq:f-form}):
\begin{equation*}
(\alpha^{a+b} - \alpha^{ab})\beta\alpha^c = 
\beta\alpha^a(\alpha^{b+c} - \alpha^{bc})
\end{equation*}
Dividing by $\beta\alpha^c$ (assuming $\beta \neq 0$): 
$\alpha^{a+b} - \alpha^{ab} = \alpha^{a+b} - \alpha^{a+bc-c}$, 
i.e., $\alpha^{ab} = \alpha^{a + bc - c}$, i.e., $ab = a + bc - c$ 
for all $a, b, c$. Setting $(a,b,c) = (2,3,4)$: $6 \neq 2 + 8 = 10$. 
Contradiction, so $\beta = 0$, hence $f = 0$ and $d^1 f = 0$.
\end{proof}

\begin{theorem}[Structure of the regular cohomology]
\label{thm:H2}
Because the complex is curved ($d^2 \neq 0$), ordinary cohomology is
not defined; we work with the \emph{regular cohomology}
$H^2_\alpha := Z^2_{\mathrm{reg}} / (Z^2_{\mathrm{reg}} \cap B^2)$.
Then $H^2_\alpha \cong \RR$, generated by $[\delta_\alpha]$.
\end{theorem}

\begin{proof}
$Z^2_{\mathrm{reg}} \cong \RR$ (Corollary~\ref{cor:Z2}), 
$Z^2_{\mathrm{reg}} \cap B^2 = \{0\}$ 
(Theorem~\ref{thm:B2-cap-Z2}), so the quotient is $\RR$.
\end{proof}

% ===================================================================
\section{Fibered structure over the parameter space}
\label{sec:fibered}
% ===================================================================

\begin{theorem}[Incompatibility of different bases]
\label{thm:fibered}
For $\alpha \neq \beta$: the cocycle $\delta_\beta$ does NOT satisfy 
the $\alpha$-cocycle condition.
\end{theorem}

\begin{proof}
The $\alpha$-cocycle condition for $g$ reads 
$\alpha^a g(b,c) - g(ab,c) + g(a,bc) - g(a,b)\alpha^c = 0$. With 
$g = \delta_\beta$, $g(a,b) = \beta^{a+b} - \beta^{ab}$:

The term $\alpha^a \beta^{b+c}$ appears in the expansion but cannot 
cancel: it would require another term of the form $c \cdot 
\alpha^a \beta^{b+c}$, but all other terms involve 
$\beta^{ab+c}$, $\beta^{a+bc}$, or $\beta^{a+b}\alpha^c$, none of 
which equals $\alpha^a\beta^{b+c}$ for generic $a, b, c$ when 
$\alpha \neq \beta$.

Numerically: $\alpha = 2$, $\beta = 3$, $(a,b,c) = (1,2,3)$, with
$g = \delta_\beta$ so that $g(x,y) = 3^{x+y} - 3^{xy}$:
\begin{align*}
d^2 g(1,2,3) &= 2^1(3^5 - 3^6) - (3^5 - 3^6) + (3^7 - 3^6)
- (3^3 - 3^2)2^3 \\
&= 2(243 - 729) - (243 - 729) + (2187 - 729) - (27 - 9)\cdot 8 \\
&= -972 + 486 + 1458 - 144 = 828 \neq 0.
\end{align*}
\end{proof}

\begin{corollary}[Fibered cohomology]
\label{cor:fibered}
The assignment $\alpha \mapsto H^2_\alpha \cong \RR$ defines a family
of one-dimensional spaces over the parameter set $(0,1) \cup
(1,\infty)$, each generated by $[\delta_\alpha]$. By
Theorem~\ref{thm:fibered} the fibers are \emph{incompatible}: for
$\alpha \neq \beta$ the generator $\delta_\beta$ of one fiber is not
even a cocycle for a different base $\alpha$, so there is no canonical
identification between distinct fibers. Endowing this family with a
genuine vector-bundle structure (topology, connection, holonomy) is
Open Problem~\ref{prob:geometric}.
\end{corollary}

% ===================================================================
\section{Connection to curved complexes}
\label{sec:curved}
% ===================================================================

\begin{definition}[{Curved cochain complex~\cite{Positselski2011}}]
\label{def:curved}
A \textbf{curved cochain complex} $(C^\bullet, d, \omega)$ consists of 
a graded module $C^\bullet$, a degree-1 map $d$, and a 
\textbf{curvature element} $\omega \in C^2$ such that $d^2 = 
[\omega, \cdot]$ (the graded commutator with $\omega$). When 
$\omega = 0$, one recovers a standard cochain complex.
\end{definition}

\begin{proposition}
\label{thm:curved}
With the differentials of Definition~\ref{def:complex} and the cup
product of Definition~\ref{def:cup}, the curvature relation
$d^2 \circ d^1 = [\delta_\varphi, \,\cdot\,]$ holds on $C^1$, and
$d^1 \circ d^0 = [\delta_\varphi, \,\cdot\,] = 0$ on $C^0$. Thus
$\delta_\varphi$ acts as a curvature element: the second-order
obstruction is the cup-commutator with $\delta_\varphi$, exactly as in
a curved cochain complex with curvature $\omega = \delta_\varphi$.
\end{proposition}

\begin{proof}
The relation on $C^1$ is Theorem~\ref{thm:non-standard}: $(d^2 \circ
d^1)(f) = \delta \smile f - f \smile \delta = [\delta, f]$. On $C^0$,
$d^0 = 0$, while scalars commute with $\delta$ under $\smile$, so
$[\delta, \,\cdot\,] = 0 = d^1 \circ d^0$.
\end{proof}

\begin{corollary}[Bianchi identity for the curvature]
\label{cor:bianchi}
The curvature element is $d^2$-closed: $d^2\delta_\varphi = 0$. Indeed, by
Definition~\ref{def:complex},
\[
(d^2\delta)(a,b,c) = \varphi(a)\delta(b,c) - \delta(ab,c) + \delta(a,bc) - \delta(a,b)\varphi(c),
\]
and this vanishes identically: it is precisely the twisted cocycle condition of
Theorem~\ref{thm:cocycle}. Thus, although $d^2$ is not a differential
($d^2\circ d^1 \neq 0$), the curvature $\delta_\varphi$ is itself a genuine
$2$-cocycle, $\delta_\varphi \in Z^2$ --- the Bianchi identity $d\omega = 0$ of a
curved complex holds in this degree. Equivalently: the first main theorem
(twisted cocycle) and the curved-complex structure are two readings of the same
identity.
\end{corollary}

\begin{remark}[What is \emph{not} claimed]
\label{rem:not-claimed}
Proposition~\ref{thm:curved} verifies the curvature relation only in
the degrees where the differentials are defined. Promoting
$(C^\bullet, \smile, d, \delta)$ to a \emph{full} curved differential
graded algebra---defining the higher differentials $d^{\ge 3}$ and
checking that $d$ is a $\smile$-derivation---is not established here (the
Bianchi identity $d^2\delta = 0$ itself \emph{does} hold, by
Corollary~\ref{cor:bianchi}); see Open Problem~\ref{prob:positselski}.
\end{remark}

\begin{remark}[Arithmetic genesis of a curved complex]
\label{rem:genesis}
Curved complexes in the literature arise from algebraic geometry 
(matrix factorizations~\cite{Eisenbud1980}), string theory 
(Landau--Ginzburg models~\cite{Orlov2004}), and Koszul 
duality~\cite{Positselski2011}. The curvature elements are typically 
algebraic or geometric objects (central elements, potentials, Massey 
products). The appearance of a curved complex from the elementary 
arithmetic operation $n \mapsto \alpha^n$ is new: the curvature 
$\delta(a,b) = \alpha^{a+b} - \alpha^{ab}$ has a purely number-theoretic 
origin.
\end{remark}

\begin{remark}[Cohomology of curved complexes]
\label{rem:curved-cohomology}
In a curved complex, $B^2 \not\subset Z^2$ in general (a coboundary is 
not automatically a cocycle). Theorem~\ref{thm:B2-cap-Z2} is a specific 
instance: the only coboundary in $Z^2_{\mathrm{reg}}$ is the trivial 
one. This is characteristic of ``non-degenerate'' curvature --- the 
curvature $\delta$ is genuinely obstruction-theoretic, not an artifact.
\end{remark}

% ===================================================================
\section{Computational verification}
\label{sec:verification}
% ===================================================================

All theorems have been verified by systematic computation (Python, 
exact arithmetic via \texttt{sympy} for symbolic, IEEE 754 double 
precision for numerical).

\begin{center}
\begin{tabular}{lrrl}
\toprule
\textbf{Result} & \textbf{Test cases} & \textbf{Failures} & 
\textbf{Method} \\
\midrule
Twisted cocycle~\eqref{eq:cocycle} & 512 triples, 5 bases & 0 & 
Exact \\
$d^2 d^1 = \delta f - f\delta$~\eqref{eq:d2d1} & 4 functions, 
30 triples & 0 & Exact \\
$[\delta_\alpha] \neq 0$ & 7 bases $\alpha$ & All: $\beta = 1/2 
\neq 1$ & Symbolic \\
Regular cocycle $\Leftrightarrow$ $\mu + \nu = 0$ & 6 pairs, 
5 triples & 0 & Exact \\
$Z^2_{\mathrm{reg}} \cap B^2 = \{0\}$ & $(2,3,4)$ test & 
$6 \neq 10$ & Exact \\
Incompatibility $\alpha \neq \beta$ & 3 pairs, 5 triples & All 
$\neq 0$ & Numerical \\
\bottomrule
\end{tabular}
\end{center}

The Python script (available as supplementary material) performs all 
verifications in under 2 seconds.

% ===================================================================
\section{Open problems}
\label{sec:open}
% ===================================================================

\begin{problem}[Full $H^2$]
\label{prob:full-H2}
Compute $H^2_\alpha$ without the exponential regularity restriction. 
Is every 2-cocycle cohomologous to a regular one?
\end{problem}

\begin{problem}[Higher cohomology]
\label{prob:higher}
Compute $H^n$ for $n \geq 3$. Does the sequence of obstructions 
$d^{n+1} \circ d^n$ define higher curvature terms (an 
$A_\infty$-structure)?
\end{problem}

\begin{problem}[Positselski axioms]
\label{prob:positselski}
Does $(C^\bullet, d, \delta)$ satisfy the axioms of a curved 
$A_\infty$-algebra in the sense of~\cite{Positselski2011}?
\end{problem}

\begin{problem}[Non-exponential transfers]
\label{prob:non-exp}
Compute $H^2$ for $\varphi(n) = p_n$ (the $n$-th prime), 
$\varphi(n) = \lfloor n\sqrt{2} \rfloor$ (Beatty sequence), or 
$\varphi(n) = n^k$ (power transfer). What replaces exponential 
regularity?
\end{problem}

\begin{problem}[Geometric interpretation]
\label{prob:geometric}
Formalize $\{H^2_\alpha\}_\alpha$ as a vector bundle. Compute its 
connection, curvature, and holonomy. Is the bundle trivializable?
\end{problem}

% ===================================================================
\section{Concluding remarks}
\label{sec:conclusion}
% ===================================================================

The defect cocycle of an arithmetic transfer exhibits a structure that 
does not fit into standard cohomological frameworks. The five key 
features are:

\begin{enumerate}
\item The cocycle is \textit{twisted}: the internal error is amplified 
exponentially by the external context.
\item The cochain complex is \textit{curved}: $d^2 \circ d^1 = 
[\delta, \cdot] \neq 0$.
\item The curvature is \textit{non-trivial}: $[\delta] \neq 0$ in 
$H^2 \cong \RR$, proved by surdetermination.
\item The parameter space is \textit{fibered}: each base $\alpha$ 
defines its own cohomology, with incompatible fibers.
\item The construction is \textit{arithmetic}: the curvature arises 
from the elementary map $n \mapsto \alpha^n$, not from geometry or 
algebra.
\end{enumerate}

The formula $d^2 \circ d^1 = [\delta, \cdot]$ places the construction 
in the theory of curved complexes~\cite{Positselski2011}, but with a 
new feature: the curvature element has a number-theoretic origin and 
carries computational information (the growth of $\delta$ controls the 
invertibility of the transfer; see~\cite{Nembe2026AGT4}).

\subsection*{Acknowledgments}

The author thanks the National Institute of Management Sciences 
(Libreville) for institutional support.

\begin{thebibliography}{20}

\bibitem{Gerstenhaber1964} M. Gerstenhaber, On the deformation of 
rings and algebras, \textit{Ann. Math.}, 79 (1964), 59--103.

\bibitem{Hochschild1945} G. Hochschild, On the cohomology groups of 
an associative algebra, \textit{Ann. Math.}, 46 (1945), 58--67.

\bibitem{Weibel1994} C. Weibel, \textit{An Introduction to 
Homological Algebra}, Cambridge University Press, 1994.

\bibitem{Loday1998} J.-L. Loday, \textit{Cyclic Homology}, 2nd ed., 
Springer, 1998.

\bibitem{Brown1982} K. Brown, \textit{Cohomology of Groups}, 
Springer, 1982.

\bibitem{MacLane1963} S. Mac Lane, \textit{Homology}, Springer, 1963.

\bibitem{Positselski2011} L. Positselski, \textit{Two Kinds of 
Derived Categories, Koszul Duality, and Comodule-Contramodule 
Correspondence}, Memoirs AMS, vol. 212, no. 996, 2011.

\bibitem{KellerLowenNicolas2009} B. Keller, W. Lowen, P. Nicol\'as, 
On the (non)vanishing of some ``derived'' categories of curved dg 
algebras, \textit{J. Pure Appl. Algebra}, 214 (2010), 1271--1284.

\bibitem{Eisenbud1980} D. Eisenbud, Homological algebra on a complete 
intersection, \textit{Trans. AMS}, 260 (1980), 35--64.

\bibitem{Orlov2004} D. Orlov, Triangulated categories of 
singularities and D-branes in Landau--Ginzburg models, \textit{Proc. 
Steklov Inst. Math.}, 246 (2004), 227--248.

\bibitem{BellierDrummond2020} J. Bellier-Mill\`es, G. C. 
Drummond-Cole, Homotopy theory of curved operads and curved algebras, 
\textit{arXiv:2007.03004}, 2020.

\bibitem{Neukirch1999} J. Neukirch, \textit{Algebraic Number Theory}, 
Springer, 1999.

\bibitem{Nembe2026AGT4} J. Nemb\'e, Defect growth of arithmetic 
transfers as a necessary condition for one-wayness, 2026.

\end{thebibliography}

\end{document}
